\begin{afterword}
\summaryhead

The post begins with an anecdote. A friend who read E.~T. Jaynes's book on probability said Jaynes was like a thousand-year-old vampire, and did not get the same sense from the author. The author agrees: Jaynes showed a level of mastery, of ``formidability'', that the author aspires to and has not reached in any field. The author suggests that people cannot easily tell apart levels far above their own.

Asked who seemed more natively intelligent than the author, Marcello Herreshoff named John Conway. The author considers three explanations and accepts the most painful one as possible: being ``strictly dumber'' than Conway. The evidence offered for accepting it is specific effort, such as blogging so that a younger mind may overtake the author. Only such effort, the author says, counts as real humility; modest words are cheap. The post ends by refusing the consolation that someone less brilliant can still do important work in the field.

\responsehead

This is a self-assessment, and it keeps its hedges. The three explanations of the Conway judgment are offered as possibilities, the rule about levels is offered as an observation, and its one number is called ``a cool-sounding wild guess''. There is little here for a critic to correct, and the notes mostly add information.

The claims about other people check out, with one qualification. Jaynes's book does line up paradoxes against Bayesian probability, in a chapter of that name. The image of Jaynes ``blasting them to pieces'' is the author's impression, and it should be read with the record: the authors of the marginalization paradox published a critique of Jaynes's treatment in 1996, and a later paper that sides with Jaynes on the conceptual point still judged his technical analysis ``not entirely successful''. On the post's own caveat, that a book shows its author's best side: Jaynes died before finishing this one, and it was edited from an incomplete manuscript. Conway's range across many fields of mathematics is as the post says.

The rule that people cannot tell apart levels far above their own has support in research. Kruger and Dunning (1999) found that the weakest performers on their tests were worse than the strongest at recognizing competence in others. Their study compares the bottom with the top, not each person with those far above, so it supports the post's rule only in part.

The account of humility is the author's own definition from ``The Proper Use of Humility'' (2006), specific actions taken against one's own errors; our notes on that post compare it with the dictionary sense.

The last line is asserted without explanation: ``that's not how it works.'' The reader of this post alone cannot tell why lesser brilliance should rule out important work. The reason comes nine days later, in ``My Bayesian Enlightenment'': the author thought Friendly AI needed ``something like the Jaynes-level'' of mastery to attempt safely.

In short: a candid, well-hedged self-assessment whose facts about others hold up, whose picture of Jaynes should be read beside the disputes over his treatment of the paradoxes, and whose closing claim depends on a reason given only in a later post.
\end{afterword}
